% Part:sets-functions-relations
% Chapter: sets
% Section: schroder-bernstein

\documentclass[../../../include/open-logic-section]{subfiles}

\begin{document}

\olfileid{sfr}{siz}{sb}

\olsection{د اندازې مفکوره او شروډر–برنشتاين}

\begin{explain}
دا يوه شهودي مفکوره ده: که $A$ له $B$ څخه لوے نۀ وي او $B$ له
$A$ څخه لوے نۀ وي، نو $A$ او $B$ همشمېره دي۔ رښتيا دا چې که
دا مفکوره \emph{ناسمه} وے، موږ به په ډېره ګرانه توجيه کړے شوے
چې زموږ تعريف شوې د همشمېر والي مفکوره د سټونو تر منځ د «اندازو»
له پرتله کولو سره څه اړيکه لري!{} خو له نېکه مرغه، شهودي مفکوره
سمه ده۔ دا د شروډر–برنشتاين په قضيې توجيه کېږي۔
\end{explain}

\begin{thm}[شروډر–برنشتاين]
	\ollabel{thm:schroder-bernstein}
	که $\cardle{A}{B}$ او $\cardle{B}{A}$ وي،
	نو $\cardeq{A}{B}$۔
\end{thm}

\begin{explain}
په بله وينا، که له $A$ څخه~$B$ ته يوه !!a{injection}، او له
$B$ څخه~$A$ ته يوه !!a{injection} وي، نو له $A$ څخه~$B$ ته
يوه !!a{bijection} شته۔

خو د دې پايلې ثبوت په رښتيا خورا \emph{ګران} دے۔ په حقيقت کښې،
که څه هم کانتور دا پايله بيان کړه، نورو ثابته کړه۔\footnote{د
تاريخ په اړه د نورو معلوماتو دپاره، د بېلګې په ډول
\citet[pp.~165--6]{Potter2004} وګورئ۔}
\oliflabeldef{sfr:cardinals:card-sb:sec}{موږ به يوازې په
\olref[sfr][cardinals][card-sb]{sec} کښې د شروډر–برنشتاين د
\emph{ثابتولو} مقام ته ورسېږو۔}{}%
اوس مهال بايد دا په باور ومنئ۔

له نېکه مرغه، شروډر–برنشتاين \emph{سمه} ده او دا فکر مو تاييدوي
چې زموږ تعريف شوې اړيکې، يعنې $\cardeq{A}{B}$ او
$\cardle{A}{B}$، له «اندازې» سره څه اړيکه لري۔ پر دې سربېره، د
شروډر–برنشتاين قضيه ډېره \emph{ګټوره} ده۔ د دوو همشمېره سټونو
تر منځ د !!a{bijection} تصورول ګرانېدلای شي۔ د شروډر–برنشتاين
قضيه اجازه راکوي چې پرتله په حالتونو ووېشو، څو يوازې له لومړي
څخه دويم ته د !!a{injection}، او بيا برعکس، تصور وکړو۔
\end{explain}
\end{document}
